\documentclass[12pt,a4paper]{article}
\usepackage[margin=25mm, headheight=15pt, headsep=10pt]{geometry}
\usepackage{fontspec}
\usepackage[table]{xcolor} % Note: table option is required for rowcolors
\usepackage{titlesec}
\usepackage{tocloft}
\usepackage{fancyhdr}
\usepackage{booktabs}
\usepackage{tcolorbox}
\tcbuselibrary{skins, breakable}
\usepackage[colorlinks=true, linkcolor=emerald, urlcolor=emerald, bookmarks=true]{hyperref}

\definecolor{emerald}{RGB}{0,102,51}
\definecolor{darkred}{RGB}{139,0,0}
\definecolor{lightgray}{RGB}{245,245,245}

\usepackage[nil]{babel}
\babelprovide[import, main]{bengali}
\babelprovide[import]{arabic}
\babelfont{rm}[Renderer=HarfBuzz,Scale=1.0]{Noto Serif Bengali}
\babelfont{sf}[Renderer=HarfBuzz]{Noto Sans Bengali}
\babelfont{tt}[Renderer=HarfBuzz]{Noto Sans Bengali}
\babelfont[arabic]{rm}[Renderer=HarfBuzz,Scale=1.1]{Noto Naskh Arabic}

% Section styling
\titleformat{\section}{\Large\bfseries\color{emerald}}{\thesection.}{0.5em}{}
\titleformat{\subsection}{\large\bfseries\color{emerald!85!black}}{\thesubsection.}{0.5em}{}
\titleformat{\subsubsection}{\normalsize\bfseries\color{emerald!70!black}}{\thesubsubsection.}{0.5em}{}
\titlespacing*{\section}{0pt}{18pt}{8pt}

% Fancy Header/Footer setup
\pagestyle{fancy}
\fancyhf{} % clear all header and footer fields
\fancyhead[L]{\color{emerald}\small\textbf{\nouppercase{\leftmark}}}
\fancyfoot[L]{\scriptsize\color{black!50}{\lat{AI}}-সংকলিত, আলেম-যাচাইকৃত নয়}
\fancyfoot[C]{\color{emerald!70!black}\thepage}
\renewcommand{\headrulewidth}{0.4pt}
\renewcommand{\headrule}{\hbox to\headwidth{\color{emerald}\leaders\hrule height \headrulewidth\hfill}}
\renewcommand{\footrulewidth}{0pt}

% Custom boxes
% 1. Ayat/Hadith Box
\newtcolorbox{ayatbox}[1][]{
    enhanced, breakable, 
    colback=emerald!5, 
    colframe=emerald, 
    boxrule=0.5pt, 
    arc=3pt, 
    left=10pt, right=10pt, top=10pt, bottom=10pt,
    #1
}

% 2. Quote/Translation Box
\newtcolorbox{quotebox}[1][]{
    enhanced, breakable,
    frame hidden,
    colback=lightgray,
    borderline west={3pt}{0pt}{emerald},
    left=15pt, right=10pt, top=10pt, bottom=10pt,
    sharp corners,
    #1
}

% 3. AI-generated notice box
\newtcolorbox{warnbox}[1][]{
    enhanced, breakable,
    colback=red!4,
    colframe=darkred,
    boxrule=0.8pt,
    arc=3pt,
    left=10pt, right=10pt, top=10pt, bottom=10pt,
    #1
}

% Commands
\newcommand{\arab}[1]{\foreignlanguage{arabic}{#1}}
\newcommand{\kib}[1]{\textcolor{emerald}{\textbf{#1}}}

% Latin (English) fallback: Noto Serif Bengali has NO Latin glyphs
\newfontfamily\latfont[Renderer=HarfBuzz]{Noto Serif}
\newcommand{\lat}[1]{{\latfont #1}}

\setlength{\parskip}{5pt}
\setlength{\parindent}{0pt}

% Fix for missing bullet in Noto Serif Bengali
\renewcommand{\labelitemi}{$\bullet$}



\begin{document}
\begin{center}{\Large\bfseries\color{emerald} পার্ট ৫ঘ — মাসআলা: জুমা, ঈদ, জানাজা, নফল, বিতর, তাহাজ্জুদ, তারাবীহ, ইস্তিখারা}\end{center}
\vspace{1em}
\section*{পার্ট ৫ঘ — মাসআলা: জুমা, ঈদ, জানাজা, নফল, বিতর, তাহাজ্জুদ, তারাবীহ, ইস্তিখারা}

\begin{warnbox}
 \textbf{গুরুত্বপূর্ণ নোটিশ (\lat{Important} \lat{Notice}):} এই কন্টেন্টটি \lat{AI} (কৃত্রিম বুদ্ধিমত্তা) দিয়ে প্রস্তুত ও সংকলিত। \lat{AI} কঠোর সূত্র-মান নীতি মেনে শুধু নির্ভরযোগ্য উৎস থেকে তথ্য সংগ্রহ করেছে — কেবল সহীহ/হাসান হাদিস ও সহীহ সনদের আসার রাখা হয়েছে, দুর্বল সনদ বাদ দেওয়া হয়েছে। তবে এটি কোনো আলেম বা মুহাদ্দিস কর্তৃক যাচাইকৃত নয়।
\end{warnbox}

\begin{quote}
\textbf{সম্পর্কিত ফাইল:} পার্ট ২-৫খ (মূল নামাজের মাসআলা), পার্ট ৫গ (সফর-নারী), পার্ট ৬ (মাযহাব-তুলনা), পার্ট ৮ (সাহাবিদের ঘটনা)।
\end{quote}

\begin{quote}
\textbf{যে প্রশ্নের উত্তর এই অংশে:} জুমার নামাজ কীভাবে? ঈদের নামাজের নিয়ম? জানাজার নামাজ কীভাবে? নফল-বিতর-তাহাজ্জুদ-তারাবীহ-ইস্তিখারার পদ্ধতি ও গুরুত্ব?
\end{quote}

\medskip\hrule\medskip

\subsection*{১. জুমার নামাজ}

\subsubsection*{১.১ ভিত্তি ও গুরুত্ব}

\textbf{সূরা আল-জুমুআহ ৬২:৯:}

\begin{center}\arab{يَـٰٓأَيُّهَا ٱلَّذِينَ ءَامَنُوٓا۟ إِذَا نُودِىَ لِلصَّلَوٰةِ مِن يَوْمِ ٱلْجُمُعَةِ فَٱسْعَوْا۟ إِلَىٰ ذِكْرِ ٱللَّهِ وَذَرُوا۟ ٱلْبَيْعَ ۚ ذَٰلِكُمْ خَيْرٌ لَّكُمْ إِن كُنتُمْ تَعْلَمُونَ}\end{center}

\textbf{উচ্চারণ:} ইয়া- আয়্যুহাল্লাযীনা আ-মানূ, ইযা- নূদিয়া লিস-সলাতি মিন ইয়াওমিল-জুমু'আতি ফাস'আও ইলা- যিকরিল্লা-হি ওয়া যারুল-বাই'। যা-লিকুম খাইরুল লাকুম ইন কুনতুম তা'লামূন।

\textbf{বাংলা অর্থ (\lat{pure}):}
\textbf{\lat{“}হে ঈমানদারগণ! যখন জুমুআর দিনে নামাজের জন্য আহ্বান করা হয়, তখন তোমরা আল্লাহর স্মরণের দিকে ধাবিত হও এবং বেচা-কেনা (ব্যবসা) ছেড়ে দাও। এটাই তোমাদের জন্য উত্তম, যদি তোমরা জানতে।\lat{”}}

\textbf{সহীহ হাদিস (সহীহ মুসলিম ৮৫৭):}

\begin{quote}
আবু হুরাইরা (রাঃ) থেকে: নবীজি (সা.) বলেন — \textbf{\lat{“}যে জুমার দিন ভালোভাবে গোসল করে, (পরিচ্ছন্ন হয়ে) মসজিদে আসে, ইমামের কাছাকাছি বসে, চুপ করে খুতবা শোনে — তার (মিনি থেকে মিনি) আগের জুমা পর্যন্ত গুনাহ মাফ করা হয়, সাথে আরও তিন দিন।\lat{”}}
\end{quote}

\subsubsection*{১.২ জুমার রাকাত ও খুতবা}

\begin{itemize}
\item \textbf{জুমার নামাজ ২ রাকাত} — জামাতে, খুতবার পর (সহীহ বুখারি ৯২৯: নবীজির (সা.) জুমার নামাজ ২ রাকাত)।
\item \textbf{খুতবা:} জুমার খুতবা \textbf{ফরজ-পূর্বশর্ত} (সব মাযহাব); খুতবা শোনা মুক্তাদিদের জন্য প্রয়োজন; খুতবার সময় কথা বলা \textbf{মাকরূহ/নিষিদ্ধ} (সহীহ বুখারি ৯৩৪: খুতবার সময় কথা বললে নামাজের সওয়াব কমে)।
\item \textbf{জুমার বিকল্প:} জুমা পেলে \textbf{যোহর পড়ার প্রয়োজন নেই} (জুমাই যোহরের স্থলাভিষিক্ত)। জুমা না পেলে/মিস হলে \textbf{যোহর ৪ রাকাত} পড়বে।
\end{itemize}
\subsubsection*{১.৩ জুমার শর্ত (সংক্ষেপ)}

\begin{itemize}
\item জুমা \textbf{ওয়াজিব} — পুরুষ মুকিমের জন্য (সব মাযহাব; হানাফিতে ওয়াজিব, হাম্বলিতে ফরজ আইন)।
\item নারী, মুসাফির, রোগী — জুমা ফরজ নয়; পড়লে জুমা জায়েয (মাযহাবগত সামান্য পার্থক্য — হানাফিতে নারীর জুমা মাকরূহ)।
\item জুমার সময়: \textbf{যোহরের সময়ে} (জুমার দিন যোহরের বদলে)।
\end{itemize}
\medskip\hrule\medskip

\subsection*{২. ঈদের নামাজ}

\subsubsection*{২.১ দুই ঈদ — ফিতর ও আজহা}

\textbf{সহীহ হাদিস (সহীহ বুখারি ৯৫২):}

\begin{quote}
আয়েশা (রাঃ) থেকে: ঈদুল ফিতরের দিন নবীজি (সা.) তাকবির দিয়ে ঈদের নামাজ পড়ান — ২ রাকাত, তাকবির সহ।
\end{quote}

\textbf{সহীহ হাদিস (সহীহ বুখারি ৯৫৬-৯৫৭):}

\begin{quote}
ঈদুল আজহায় নবীজি (সা.) ঈদগাহে গিয়ে নামাজ পড়ান — তারপর খুতবা দেন।
\end{quote}

\subsubsection*{২.২ ঈদের পদ্ধতি (সংক্ষেপ)}

\begin{enumerate}
\item ঈদগাহ/মসজিদে — \textbf{আযান-ইকামত ছাড়া} ২ রাকাত (সহীহ মুসলিম ৮৮৫: জাবির — আযান-ইকামত ছাড়া)।
\item রাকাতে \textbf{অতিরিক্ত তাকবির:}
\end{enumerate}
\begin{itemize}
\item \textbf{হানাফি:} ৩ বার তাকবির (রুকুতে যাওয়ার আগে), প্রথম রাকাতে সানা-এর আগে ৩ তাকবির, দ্বিতীয় রাকাতে রুকুতে যাওয়ার আগে ৩ তাকবির।
\item \textbf{জমহুর (শাফেয়ি-মালেকি-হাম্বলি):} প্রথম রাকাতে সানা-এর পরে ৬-৭ তাকবির, দ্বিতীয় রাকাতে রুকুর আগে ৫ তাকবির (সংখ্যায় পার্থক্য আছে)।
\item (সঠিক সংখ্যা-পার্থক্য মাযহাবগত — দলিল দুই রকম; নিয়ম: ইমামের অনুসরণ)
\end{itemize}
\begin{enumerate}
\item \textbf{খুতবা:} ঈদের খুতবা \textbf{নামাজের পরে} (জুমার খুতবা আগে — পার্থক্য লক্ষ করুন)।
\item \textbf{ঈদ-প্রস্তুতি:} ঈদের দিন সকালে খেয়ে ঈদগাহে যাওয়া (ফিতরে), পথে তাকবির (সহীহ-হাসান বর্ণনা — সূরা বাকারাহ ২:১৮৫-এর আলোকে)।
\end{enumerate}
\subsubsection*{২.৩ ঈদের দিনের সুন্নত-আদব}

\begin{itemize}
\item ঈদুল ফিতরে নামাজের \textbf{আগে} মিষ্টি কিছু খাওয়া (সহীহ বুখারি ৯৫৩)।
\item ঈদগাহে \textbf{ভিন্ন পথে} যাওয়া-আসা (সহীহ বুখারি ৯৮৬)।
\item গোসল, উত্তম পোশাক, তাকবিরে মুখর হওয়া।
\end{itemize}
\medskip\hrule\medskip

\subsection*{৩. জানাজার নামাজ}

\subsubsection*{৩.১ গুরুত্ব}

\textbf{সহীহ হাদিস (সহীহ মুসলিম ৯৪৫):}

\begin{quote}
নবীজি (সা.) বলেন — \textbf{\lat{“}যে ব্যক্তি জানাজায় অংশ নেয় এমনকি দাফন পর্যন্ত, সে কীরাত (দুই পাহাড় পরিমাণ) সওয়াব পায়; যে শুধু জানাজার নামাজে থাকে, সে এক কীরাত পায়।\lat{”}} (কীরাত = উহুদ পাহাড় পরিমাণ সওয়াব)।
\end{quote}

\subsubsection*{৩.২ জানাজার নামাজের পদ্ধতি}

\begin{itemize}
\item \textbf{৪ তাকবির} (সব মাযহাবের প্রসিদ্ধ — সহীহ বুখারি ১২৯৭-এ ৪ তাকবিরের বর্ণনা) — \textbf{রুকু-সিজদা নেই}, সবটাই দাঁড়িয়ে।
\item \textbf{ধাপ:}
\end{itemize}
\begin{enumerate}
\item নিয়ত (হৃদয়ে)।
\item \textbf{১ম তাকবির:} সানা (তাকবিরের পর) — অতঃপর আউযুবিল্লাহ + ফাতিহা (জমহুর; হানাফিতে ফাতিহা তিলাওয়াত)।
\item \textbf{২য় তাকবির:} দুরুদ (নবীজির (সা.) উপর)।
\item \textbf{৩য় তাকবির:} মৃত ব্যক্তির জন্য দুআ (অতএব আল-ফাতিহা দুআ-পদ্ধতিতে — \lat{“}আল্লাহুম্মাগফির লিহাইয়্যিনা...\lat{”} — হাদিসে মুসলিম ৯৬৩)।
\item \textbf{৪র্থ তাকবির:} সালাম (ডানে) — (সালাম দুই দিকে মাযহাবভেদে)।
\end{enumerate}
\begin{itemize}
\item \textbf{শর্ত:} জানাজার নামাজের জন্য \textbf{পবিত্রতা (ওযু/গোসল)} শর্ত (হানাফি-জমহুর; মুসলিম-সহিহ হাদিসে নবীজি জানাজার নামাজে ওযু করতেন)। মুসলিম ৫৭৬-এ নবীজির ওযু করে জানাজার নামাজের বর্ণনা।
\item \textbf{সময়:} মৃত ব্যক্তিকে গোসল-কাফন করার পর; দাফনের আগে/পরে (পরবর্তীতে হলে পার্থক্য আছে)।
\end{itemize}
\subsubsection*{৩.৩ মহিলা-শিশু জানাজা}

\begin{itemize}
\item \textbf{মহিলা মৃত:} মাযমানে নারীরা জানাজার নামাজ পড়ে (সহীহ বুখারি ১২৮৪-তে নবীজি নাজাশীর জানাজা পড়েছেন; নারী-মৃতের ক্ষেত্রেও নিয়ম একই)।
\item \textbf{শিশু:} জানাজার নামাজ আছে (সহীহ বুখারি ১২০৫: নবীজি শিশু-মৃতের জানাজা পড়েছেন)।
\end{itemize}
\medskip\hrule\medskip

\subsection*{৪. নফল নামাজ}

\subsubsection*{৪.১ নফল (সুন্নত-রাকাত) — নির্ধারিত}

\begin{table}[h]\centering\small
\begin{tabular}{lll}
নামাজ & রাকাত & সময় \\
\hline
ফজরের সুন্নত & ২ (মুআক্কাদা) & ফজরের ফরজের আগে \\
যোহরের সুন্নত & ৪ আগে (মুআক্কাদা) + ২ পরে & যোহরের ফরজের আগে-পরে \\
আসরের সুন্নত & ২-৪ (গায়রে মুআক্কাদা; নবীজি ২ রাকাত পড়তেন — বুখারি ৯৩৭) & আসরের আগে \\
মাগরিবের সুন্নত & ২ (মুআক্কাদা) & মাগরিবের ফরজের পরে \\
ইশার সুন্নত & ২ (মুআক্কাদা) & ইশার ফরজের পরে \\
\textbf{বিতর} & ৩ (হানাফি: ওয়াজিব; জমহুর: সুন্নতে মুআক্কাদা) & ইশার পরে — রাতের শেষাংশ পর্যন্ত \\
\hline
\end{tabular}
\end{table}

\textbf{ভিত্তি হাদিস:} \lat{“}যে ব্যক্তি (নিশ্চিত) ১২ রাকাত নফল প্রতিদিন পড়বে — যোহরের আগে ৪ ও পরে ২, মাগরিবের পরে ২, ইশার পরে ২, ফজরের আগে ২ — আল্লাহ তার জন্য জান্নাতে একটি ঘর বানাবেন\lat{”} (সুনানে তিরমিজি ৪১৪ — হাসান সহীহ)।

\subsubsection*{৪.২ বিতর}

\textbf{সহীহ হাদিস (সহীহ বুখারি ৯৯০; সহীহ মুসলিম ৭৫২):}

\begin{quote}
নবীজি (সা.) বলেন — \textbf{\lat{“}বিতর তোমাদের উপর ওয়াজিব (একান্ত প্রয়োজনীয়), যারা বিতর পড়ে না তাদের (নামাজ) কষ্টকর হয়ে যায়।\lat{”}} — হানাফি এ থেকেই বিতরকে \textbf{ওয়াজিব} বলেন।
\end{quote}

\begin{itemize}
\item \textbf{পদ্ধতি:} ১ রাকাত অথবা ৩ রাকাত (সহীহ হাদিস — দুটিই প্রমাণিত); হানাফি ৩ রাকাত (বিতর) — তৃতীয় রাকাতে \textbf{কুনুত} (দুআ)।
\item \textbf{কুনুত:} হানাফি মতে বিতরের তৃতীয় রাকাতে কুনুত ওয়াজিব; জমহুরে সুন্নত। (দুআ-কুনুত — হাসান-সহীহ হাদিসে নবীজি বিতরে কুনুত পড়েছেন — সুনানে আবু দাউদ ১৪২৫)।
\end{itemize}
\subsubsection*{৪.৩ তাহাজ্জুদ (রাতের নামাজ)}

\textbf{সূরা আল-মুজ্জাম্মিল ৭৩:২-৩ (মূল আয়াত):}

\begin{center}\arab{قُمِ ٱلَّيْلَ إِلَّا قَلِيلًا ۝ نِّصْفَهُۥٓ أَوِ ٱنقُصْ مِنْهُ قَلِيلًا}\end{center}

\textbf{উচ্চারণ:} কুমিল-লাইলা ইল্লা- ক্বলীলা — নিসফাহূ আও আনকুস মিনহু ক্বলীলা।

\textbf{বাংলা অর্থ (\lat{pure}):}
\textbf{\lat{“}রাত্রির (নামাজে) দাঁড়াও — সামান্য অংশ ছাড়া; অর্ধেক রাত, অথবা তার চেয়ে একটু কম।\lat{”}}

\textbf{সহীহ হাদিস (সহীহ বুখারি ১১৩০; সহীহ মুসলিম ৭৪১):}

\begin{quote}
আয়েশা (রাঃ) থেকে: নবীজি (সা.) রাতের নামাজ পড়তেন — \textbf{যতক্ষণ পা ফুলে যেত}; বলা হলো — আপনার অতীত-ভবিষ্যৎ গুনাহ মাফ, তবু এত কষ্ট? নবীজি (সা.) বললেন: \textbf{\lat{“}আমি কি কৃতজ্ঞ বান্দা হব না?\lat{”}} (পূর্ণ ঘটনা পার্ট ৮-এর ফাইল ১)।
\end{quote}

\begin{itemize}
\item \textbf{সময়:} রাতের শেষাংশ (তাহাজ্জুদ) — ফজরের আগে।
\item \textbf{রাকাত:} ২, ৪, ৬, ৮... (জোড়-জোড়); সর্বোচ্চ কোন নির্দিষ্ট সীমা নবীজির আমল থেকে নির্ধারিত নয় — শেষে বিতর।
\end{itemize}
\subsubsection*{৪.৪ তারাবীহ (রমজানের রাতের নামাজ)}

\begin{itemize}
\item \textbf{ভিত্তি:} নবীজি (সা.) রমজানে রাতের নামাজ পড়েছেন — জামাতে ২-৩ রাত, পরে (জামাতে আরও পড়লে ফরজ হয়ে যাওয়ার ভয়ে) ছেড়ে দিয়েছেন — (সহীহ বুখারি ২০১২; সহীহ মুসলিম ৭৬১)।
\item \textbf{ফজিলত হাদিস (সহীহ বুখারি ৩৭; সহীহ মুসলিম ৭৫৯):} \textbf{\lat{“}যে ব্যক্তি ঈমান ও সওয়াবের আশায় রমজানের রাতে (তারাবীহ) দাঁড়ায়, তার আগের গুনাহ মাফ।\lat{”}}
\item \textbf{রাকাত সংখ্যা — বাস্তব ইখতিলাফ:}
\item \textbf{হানাফি-শাফেয়ি-হাম্বলি (প্রসিদ্ধ):} ২০ রাকাত (সহীহ সনদে উমার (রাঃ)-এর যুগে ২০ রাকাত জামাতের প্রমাণ — মালিকের মাওয়াত্তা, সহীহ সনদ)।
\item \textbf{মালেকি:} ৩৬ রাকাত (মদিনার আমল-ভিত্তিক) — কিছু বর্ণনায় ২০ও।
\item \textbf{সংক্ষেপে:} ২০ রাকাতই ব্যাপকভাবে প্রমাণিত; তবে রাতের শেষে \textbf{বিতর} মিলিয়ে জামাতের পূর্ণতা রক্ষা — আরও বিস্তারিত পার্ট ৬।
\end{itemize}
\subsubsection*{৪.৫ ইস্তিখারা}

\textbf{সহীহ হাদিস (সহীহ বুখারি ১১৬২; সহীহ মুসলিম ২৬০৮):}

\begin{quote}
নবীজি (সা.) বলেন — \textbf{\lat{“}তোমাদের কেউ কোনো গুরুত্বপূর্ণ কাজের সিদ্ধান্ত নিতে চাইলে, সে যেন ২ রাকাত নফল পড়ে তারপর এই দুআ করে: "আল্লাহুম্মা ইন্নি আসতাখীরুকা..."\lat{”}} — ইস্তিখারা ২ রাকাত + বিশেষ দুআ; \textbf{আত্মবিশ্বাস ও আল্লাহর উপর ভরসার} শিক্ষা।
\end{quote}

\medskip\hrule\medskip

\subsection*{৫. প্রশ্নোত্তর}

\begin{table}[h]\centering\small
\begin{tabular}{ll}
প্রশ্ন & উত্তর \\
\hline
জুমা ২ রাকাত কেন? & জুমার ফরজ ২ রাকাত (খুতবা + নামাজ) — যোহরের বদলে \\
ঈদের নামাজে অতিরিক্ত তাকবির কয়টা? & মাযহাবভেদে ৩/৬-৭/৫ — ইমামের অনুসরণ করবেন \\
জানাজার নামাজ কয় তাকবির? & ৪ তাকবির, রুকু-সিজদা নেই \\
বিতর কয় রাকাত? & ১ বা ৩ (হানাফি: ৩ + কুনুত) \\
তারাবীহ কত রাকাত? & ২০ (প্রসিদ্ধ; মালেকি ৩৬) — সবই দলিলভিত্তিক বৈধ \\
ইস্তিখারা কীভাবে? & ২ রাকাত নফল + হাদিসের দুআ \\
\hline
\end{tabular}
\end{table}

\medskip\hrule\medskip

\subsection*{৬. রেফারেন্স}

\begin{itemize}
\item কুরআন: আল-জুমুআহ ৬২:৯; আল-মুজ্জাম্মিল ৭৩:২-৩; আল-বাকারাহ ২:১৮৫
\item সহীহ বুখারি — ৩৭, ৯২৯, ৯৩৪, ৯৩৭, ৯৫২, ৯৫৩, ৯৫৬-৯৫৭, ৯৮৬, ৯৯০, ১১৩০, ১১৬২, ১২০৫, ১২৯৭, ২০১২
\item সহীহ মুসলিম — ৭৫২, ৭৫৯, ৭৬১, ৮৫৭, ৮৮৫, ৯৪৫, ৯৬৩, ২৬০৮
\item সুনানে তিরমিজি — ৪১৪ (১২ রাকাত নফল); সুনানে আবু দাউদ — ১৪২৫ (কুনুত)
\item আল-ফিকহুল ইসলামি ওয়া আদিল্লাতুহু (আল-যুহাইলি) — মাযহাব-সারণি
\end{itemize}
\begin{quote}
\textbf{দ্রষ্টব্য:} তারাবীহ-বিতরের রাকাত-সংখ্যার পূর্ণ মাযহাব-তুলনা ও দলিল পার্ট ৬-এ; ঈদের তাকবিরের দলিলভিত্তিক বিশ্লেষণও পার্ট ৬-এ।
\end{quote}

\end{document}
